\clearpage
\begin{trapbox}
\begin{center}\begin{tikzpicture}[x=.85cm,y=.85cm,>=stealth,font=\small]
\path[use as bounding box] (-7,-2) rectangle (7,3.5);
\node[font=\bfseries] at (0,3.2) {同尺度比较：动点在两端点之间，还是外侧？};
\begin{scope}[xshift=-3.5cm]
\draw[->,black!45] (-3.2,0)--(3.2,0);
\draw[blue!65!black,thick,domain=0:360,samples=120,variable=\t] plot ({2*cos(\t)},{sin(\t)});
\draw[orange!85!black,very thick] (-2,0)--(1,.866025);
\draw[teal!80!black,very thick] (2,0)--(1,.866025);
\fill (1,.866025) circle (2pt) node[above] {$P$};
\fill (-2,0) circle (2pt) node[below left=5pt] {$A_1$};
\fill (2,0) circle (2pt) node[below right=5pt] {$A_2$};
\node at (0,2.2) {$x^2/4+y^2=1$};
\node at (0,-1.65) {$-2<x_P<2,\quad k_1k_2=-1/4$};
\end{scope}
\begin{scope}[xshift=3.5cm]
\draw[->,black!45] (-3.2,0)--(3.2,0);
\foreach \s in {-1,1} {\draw[blue!65!black,thick,domain=-1.2:1.2,samples=70,variable=\t] plot ({\s*sqrt(4+4*\t*\t)},\t);}
\draw[orange!85!black,very thick] (-2,0)--(3,1.118034);
\draw[teal!80!black,very thick] (2,0)--(3,1.118034);
\fill (3,1.118034) circle (2pt) node[above] {$P$};
\fill (-2,0) circle (2pt) node[below right=3pt] {$A_1$};
\fill (2,0) circle (2pt) node[below left=3pt] {$A_2$};
\node at (0,2.2) {$x^2/4-y^2=1$};
\node at (0,-1.65) {$|x_P|>2,\quad k_1k_2=+1/4$};
\end{scope}
\end{tikzpicture}
\end{center}
图中横向端点相同；椭圆两斜率异号，双曲线两斜率同号。

\begin{description}[style=nextline,leftmargin=2.8em,itemsep=0.55em]
 \item[\textbf{\textcolor{CTrap}{易错1：误写负号（与椭圆混淆）}}]
 \textbf{错：}\(k_{PA_1}k_{PA_2}=-b^2/a^2\)。
 \textbf{正：}双曲线\(x^2/a^2-y^2/b^2=1\)对应
 \(y^2=(b^2/a^2)(x^2-a^2)\)，所以\(\boxed{k_1k_2=+b^2/a^2}\)。
 椭圆方程才会出现负号。

 \item[\textbf{\textcolor{CTrap}{易错2：忘记排除实轴端点}}]
 顶点\((\pm a,0)\)处，一条连线退化，斜率根本不定义。
 \textbf{正确处理：}\(P\in C\setminus\{A_1,A_2\}\)，
 等价于曲线上\(x\ne\pm a\)，此时\(y\ne0\)。
 逆向轨迹也必须删除这两个点。

 \item[\textbf{\textcolor{CTrap}{易错3：把两定点当成焦点}}]
 \textbf{错：}在焦点\((\pm c,0)\)处使用同样的定值。
 \textbf{正：}本结论指定\textbf{实轴端点}\((\pm a,0)\)，不是\((\pm c,0)\)；
 换定点以后分母是\(x^2-c^2\)，不能约出相同定值。

 \item[\textbf{\textcolor{CTrap}{易错4：横纵双曲线不分}}]
 \textbf{横轴型}\(x^2/a^2-y^2/b^2=1\)：\(k_1k_2=b^2/a^2=e^2-1\)。
 \textbf{纵轴型}\(y^2/a^2-x^2/b^2=1\)：\(k_1k_2=a^2/b^2=1/(e^2-1)\)。
 \textbf{原因：}斜率总是\(\Delta y/\Delta x\)，不随坐标轴方向自动取倒数。

 \item[\textbf{\textcolor{CTrap}{易错5：斜率和求坐标时漏掉交换顺序}}]
 若已算出无序斜率\(\{u,v\}=\{\alpha,\beta\}\)，
 必须检查\((u,v)=(\alpha,\beta)\)与\((\beta,\alpha)\)。
 对例3会得到两个不同点\((3,\frac52)\)与\((-3,-\frac52)\)，
 而不是一个。

 \item[\textbf{\textcolor{CTrap}{易错6：只检查韦达判别式\(\Delta\ge0\)}}]
 两条连线若斜率相同且都过\(P\)，就必须重合；此时\(A_1,A_2,P\)共线，
  \(P\)只能在实轴上，这与\(P\)是非顶点的曲线上点矛盾。
 \textbf{必须}\(u\ne v\)，所以对斜率和\(s\)、斜率积\(m>0\)，
 有\(\boxed{s^2>4m}\)。等号只对应不可能实现的“斜率相同”。

 \item[\textbf{\textcolor{CTrap}{易错7：错误声称斜率乘积为常数可判定整条双曲线}}]
 逆向判断需两斜率均存在，且\(m>0\)。
 轨迹与双曲线方程作为\textbf{斜率命题}等价时，要明确\textbf{不含顶点}。

 \item[\textbf{\textcolor{CTrap}{易错8：圆幂法在椭圆内仍作实切线}}]
 双曲线非顶点时 $|x|>a$，垂足 $H$ 在以 $A_1A_2$ 为直径的圆\textbf{外}，切线 $HT$ 存在；
 椭圆非顶点则 $|x|<a$，$H$ 在圆\textbf{内}，\textbf{不能照搬圆外点作切线的步骤}。
 这也解释了斜率积的正负区别。圆幂法最后仍要用曲线方程给出 $PH/HT=b/a$。

 \item[\textbf{\textcolor{CTrap}{易错9：把统一模型的所有椭圆都看成横长轴}}]
 固定 $A_1(-a,0),A_2(a,0)$ 时，
 $m<0$ 推出 $\frac{x^2}{a^2}+\frac{y^2}{a^2|m|}=1$。
 若 $|m|>1$，椭圆的长轴在 $y$ 方向，不能把 $A_1,A_2$ 称为\textbf{长轴端点}。
 特别地 $m=-1$ 是\textbf{圆}，轨迹原题仍须排除两端点。

 \item[\textbf{\textcolor{CTrap}{易错10：半角参数把顶点误包含进来}}]
 右支参数 $x=a\sec\theta,y=b\tan\theta$，
 必须 $-\pi/2<\theta<\pi/2$ 且 $\theta\ne0$。
 $\theta=0$ 对应 $P=A_2$，此时斜率式里有 $0/0$，不能先约分再声称顶点可用。

 \item[\textbf{\textcolor{CTrap}{易错11：将渐近线的极限当成有限点}}]
 $k_1,k_2$ 在 $P$ 趋向无穷远处会\textbf{同时趋近} $\pm b/a$，
 但有限的非顶点 $P$ 下始终 $k_1\ne k_2$。
 故给定 $m>0$ 时 $|s|>2\sqrt m$ 是严格不等式；边界\textbf{不是可取的最值}。

\end{description}
\end{trapbox}
